\documentclass[12pt]{article}
\usepackage[utf8]{inputenc}
\usepackage[T1]{fontenc}
\usepackage[margin=1in]{geometry}
\usepackage{setspace}
\usepackage{xcolor}
\usepackage[colorlinks=true, urlcolor=blue!60!black, linkcolor=blue!60!black]{hyperref}
\setlength{\parindent}{0em}
\setlength{\parskip}{0.45em}
\renewcommand{\baselinestretch}{1.0}
\begin{document}
\begin{center}
{\Large \textit{R.A.V. v. City of St. Paul}, 505 U.S. 377 (1992)\par}
\vspace{0.3em}
{\normalsize Case Brief --- One-Page Edition\par}
\end{center}

\textbf{Memory Jogger.} The cross-burning case: no content discrimination within fighting words (\href{https://supreme.justia.com/cases/federal/us/505/377/}{505 U.S. 377}).

\textbf{Facts.} A juvenile allegedly burned a cross in a Black family's yard in St.~Paul and was charged under an ordinance banning symbols that arouse anger, alarm, or resentment on the basis of race, color, creed, religion, or gender.

\textbf{Procedural History.} The Minnesota Supreme Court upheld the ordinance, limiting it to fighting words.

\textbf{Votes.} Reversed, 9--0 in judgment, 5--4 on rationale. Scalia wrote for the Court; White, Blackmun, and Stevens concurred in the judgment.

\textbf{Holding.} Even within proscribable fighting words, government may not single out disfavored topics or viewpoints; the ordinance fails strict scrutiny.

\textbf{Reasoning.} Unprotected categories may be regulated for their proscribable content, not their message, subject to narrow exceptions. The ordinance targeted five topics and, in practice, one side of a debate, and a neutral ban on all fighting words would serve the City's interest. White would have found it overbroad.

\textbf{Examples.} \emph{Same-Side:} banning only fighting words against police. \emph{Opposite-Side:} enhanced penalties for threats against poll workers. \emph{Fence-Sitter:} a topic-based campus ban.

\textbf{Critique.} Critics note the paradox of preferring broader bans; defenders praise protection of disfavored views.

\textbf{Key Quotation.} ``What they mean is that these areas of speech can, consistently with the First Amendment, be regulated because of their constitutionally proscribable content (obscenity, defamation, etc.)---not that they are categories of speech entirely invisible to the Constitution, so that they may be made the vehicles for content discrimination unrelated to their distinctively proscribable content.'' --- Justice Scalia, Opinion of the Court, 505~U.S.~at~383--384.

\end{document}
